\documentclass[10pt,a4paper]{amsart}
\usepackage[margin=19mm]{geometry}
\usepackage{amsmath,amssymb,mathtools}
\usepackage[scr=rsfs,cal=euler]{mathalfa}
\usepackage{xfrac}
\usepackage[unicode,hidelinks,bookmarks=false]{hyperref}
\newcommand{\ie}{i.e.\ }
\newcommand{\cf}{cf.\ }
\newcommand{\resp}{respectively\ }
\newcommand{\Subsection}{\subsection}
\providecommand{\nobreakdash}{\nobreak\hbox{-}}
\setlength{\emergencystretch}{2em}
\multlinegap0pt
\usepackage[shortlabels]{enumitem}
\usepackage{xcolor}
\usepackage{amssymb}
\usepackage{mathtools}
\newtheorem{lemma}{Lemma}
\newcommand{\Prb}{\mathbb P}
\newcommand{\R}{\mathbb R}
\newcommand{\Unif}{\operatorname{Unif}}
\newcommand{\dtv}{d_{\mathrm{TV}}}
\newcommand{\one}{\mathbf 1}
\newcommand{\vol}{\operatorname{vol}}
\title{Hit-and-Run Mixes as Fast as the Ball Walk}
\author{Ruizhe Zhang}
\date{Source: arXiv:2608.13487v2 (CC BY 4.0)}
\begin{document}
\maketitle

\noindent\emph{Source and licence.} Hit-and-Run Mixes as Fast as the Ball Walk. Version arXiv:2608.13487v2 (CC BY 4.0), distributed by arXiv under the Creative Commons Attribution 4.0 International licence, http://creativecommons.org/licenses/by/4.0/, as recorded in the arXiv licence metadata for this exact version and shown on the abstract page https://arxiv.org/abs/2608.13487v2. Full licence notice: \texttt{licenses/arxiv-2608.13487-CC-BY-4.0.txt}.\par\smallskip

\noindent\textit{Editorial scope.} Counted unit: the article's lemma lem:log-affine-local-overlap (source0.tex lines 3363--3372). The article's complete original proof of it is delivered with it (source0.tex lines 3374--3538, one \texttt{proof} environment of 165 lines). No mathematics is written, repaired or re-derived: the statement and the proof are the article's own lines, in the article's own notation, and no retained line is substituted or re-flowed.  Retained uncounted as the source-local context the proof needs: lines 3106--3129.  Nothing was omitted from any retained span, so the package carries no omission record.  No retained line contains a citation, so the wrapper carries no bibliography and no bibliography entry.  No retained line contains a figure environment, an image-inclusion command or drawing code, so no image is delivered and the package declares no asset dependency.  Editorial formatting only, in addition: an AMS \texttt{amsart} preamble that restates the article's own declarations and macros used by the retained lines, namely the environments \texttt{lemma} on separate counters, together with the standard packages those lines need.  The retained environments print in this document as Lemma 1 in order of appearance, whereas the article prints the same items with the numbering of its own full document.  The complete native source of the article as distributed in the arXiv e-print for this exact version is bundled unchanged as \texttt{sources/207594/source0.tex}.
\begin{lemma}[Local geometry of the intrinsic metric]
\label{lem:log-affine-intrinsic-local-geometry}
Let $L=O(\sqrt{n})$ be the Lipschitz parameter for $\rho_b$.  For every $x,y\in\operatorname{int}(K)$,
\begin{equation}\label{eq:log-affine-intrinsic-log-distortion}
 \left|\log\frac{\rho_b(y)}{\rho_b(x)}\right|
 \le L d_{\rho_b}(x,y)
\end{equation}
and
\begin{equation}\label{eq:log-affine-intrinsic-euclidean-distortion}
 \|x-y\|_2
 \le
 \rho_b(x)e^{L d_{\rho_b}(x,y)}d_{\rho_b}(x,y).
\end{equation}
Consequently, for a universal $c_*>0$, if
\begin{equation}\label{eq:log-affine-intrinsic-nearby}
 d_{\rho_b}(x,y)<\frac{c_*}{\sqrt n}
\end{equation}
then
\begin{equation}\label{eq:log-affine-intrinsic-local-geometry}
 \frac12\rho_b(x)\le\rho_b(y)\le2\rho_b(x),
 \qquad
 \|x-y\|_2\lesssim \frac{\rho_b(x)}{\sqrt n}.
\end{equation}
\end{lemma}

\begin{lemma}[Local overlap of the comparison kernel]
\label{lem:log-affine-local-overlap}
There are universal constants $c,\alpha_0>0$ such that
\begin{equation}\label{eq:log-affine-local-overlap}
 d_{\rho_b}(x,y)<\frac{c}{\sqrt n}
 \quad\Longrightarrow\quad
 \dtv\bigl({\sf P}_{\rm loc}(x,\cdot),{\sf P}_{\rm loc}(y,\cdot)\bigr)
 \le1-\alpha_0.
\end{equation}
\end{lemma}

\begin{proof}
Lemma~\ref{lem:log-affine-intrinsic-local-geometry} gives
$\rho_b(y)\in[\rho_b(x)/2,2\rho_b(x)]$ and
$\|x-y\|_2\le c\rho_b(x)/\sqrt n$.  Let
\[
 \rho_-:=\min\{\rho_b(x),\rho_b(y)\},
 \qquad
 \rho_+:=\max\{\rho_b(x),\rho_b(y)\}.
\]
Since $\rho_+/\rho_-\le2$, the interval
$[\rho_+/32,\rho_-/8]$ contains at least one dyadic number.  Choose $j$ so that
\begin{equation}\label{eq:log-affine-overlap-scale}
 \frac{\rho_+}{32}\le\delta_j\le\frac{\rho_-}{8}.
\end{equation}

Take any $x,y\in D_j$.  For $z\in \{x,y\}$, let 
\begin{align*}
    L_z:=\{w:f_b(w)\ge\beta f_b(z)\}.
\end{align*}
The commonly accepted region of the two kernels is
\[
 G:=B(x,\delta_j)\cap B(y,\delta_j)\cap D_j\cap L_x\cap L_y.
\]
Our goal is to show that both ${\sf P}_{\rm loc}(x,\cdot)$ and ${\sf P}_{\rm loc}(y,\cdot)$ have large mass on $G$.

We first bound the volume of the overlap $B(x,\delta_j)\cap B(y,\delta_j)$. Let $e:=\frac{y-x}{\|y-x\|_2}$.  We integrate along lines parallel to $e$, indexed by $z\in e^\perp$.  For $\|z\|_2\le\delta_j$, the sections of $B(x,\delta_j)$ and $B(y,\delta_j)$ on the line $z+\R e$ are intervals of the same length, translated from one another by $d$.  Hence, the length lost from their intersection is at most $d$.  Fubini's theorem
therefore gives
\[
 v_n\delta_j^n
 -\vol\bigl(B(x,\delta_j)\cap B(y,\delta_j)\bigr)\le
 d\,v_{n-1}\delta_j^{n-1}.
\]
Since $v_{n-1}/v_n\leq O(\sqrt n)$,
\[
 \frac{\vol(B(x,\delta_j)\cap B(y,\delta_j))}
 {v_n\delta_j^n}
 \ge
 1-O(\sqrt n)\frac{\|x-y\|_2}{\delta_j}\geq 1-O\left(\frac{c\rho_b(x)}{\delta_j}\right)\geq 1-O(c).
\]
For sufficiently small $c$, the normalized overlap is at least $9/10$.

We next control the losses caused by the band and density conditions. Fix $z\in\{x,y\}$ and let $V\sim\Unif(B(0,\delta_j))$.  Since
\[
 2\delta_j
 \le\frac{\rho_-}{4}
 <\rho_b(z),
\]
monotonicity of
$r\mapsto\lambda_b(z,r)$ and the definition of $\rho_b(z)$ give
\[
 \lambda_b(z,\delta_j)\ge \frac{63}{64},
 \qquad
 \lambda_b(z,2\delta_j)\ge\frac{63}{64}.
\]
The random variables $z-V$ and $z+2V$ are uniform in
$B(z,\delta_j)$ and $B(z,2\delta_j)$, respectively.  Dropping the
density condition from the definition of $\lambda_b$, we obtain
\[
 \Prb\{z-V\notin K\}\le \frac{1}{64},
 \qquad
 \Prb\{z+2V\notin K\}\le \frac{1}{64}.
\]
By the union bound, with probability at least $31/32$, both points belong to $K$.

On this event, by the concavity and non-negativity of $\rho_b$, we have
\begin{align*}
    \rho_b(z+V)\geq \frac{\rho_b(z)+\rho_b(z+2V)}{2}\geq \frac{1}{2}\rho_b(z)\,\qquad \rho_b(z)\geq \frac{\rho_b(z-V)+\rho_b(z+V)}{2}\geq \frac{1}{2}\rho_b(z+V).
\end{align*}
That is,
\[
 \frac12\rho_b(z)
 \le
 \rho_b(z+V)
 \le
 2\rho_b(z).
\]
The choice of $\delta_j$ gives, for either $z=x$ or $z=y$,
\[
 4\delta_j
 \le\frac{1}{2}\rho_b(z)\leq \rho_b(z+V)\leq 2\rho_b(z)
 \le64\delta_j.
\]
Consequently, on the preceding event,
\[
 z+V\in D_j.
\]
Since $z+V$ is uniform in $B(z,\delta_j)$, this proves
\begin{equation}\label{eq:band-retention}
 \vol(B(z,\delta_j)\setminus D_j)
 \le
 \frac{1}{32} v_n\delta_j^n,\qquad z\in \{x,y\}.
\end{equation}


Because $\delta_j<\rho_b(z)$,
\[
 \lambda_b(z,\delta_j)
 =
 \frac{\vol(B(z,\delta_j)\cap L_z)}{v_n \delta_j^n}
 \ge\frac{63}{64}.
\]
Hence
\begin{equation}\label{eq:density-retention}
 \vol(B(z,\delta_j)\setminus L_z)
 \le \frac{1}{64}v_n \delta_j^n.
\end{equation}

Combining them together, we have 
\begin{align*}
 \vol(G)
 &\ge
 \vol(B(x,\delta_j)\cap B(y,\delta_j))
 -\vol(B(x,\delta_j)\setminus D_j)
 -\vol(B(x,\delta_j)\setminus L_x)
 -\vol(B(y,\delta_j)\setminus L_y)\\
 &\ge
 \left(
  \frac9{10}-\frac{1}{32} - \frac{2}{64}
 \right)v_n\delta_j^n = \frac{67}{80}v_n\delta_j^n.
\end{align*}

For every $w\in G$, we have
\[
 x,y,w\in D_j,
 \qquad
 \|w-x\|_2,\|w-y\|_2\le\delta_j,\qquad
 f_b(w)\ge\beta f_b(x),
 \qquad
 f_b(w)\ge\beta f_b(y).
\]
Therefore, the $j$th term of the transition density satisfies
\[
 k(x,w)\ge\frac{\alpha\beta}{v_n\delta_j^n},
 \qquad
 k(y,w)\ge\frac{\alpha\beta}{v_n\delta_j^n}.
\]
Thus, the measure
\[
 d\xi(w)
 :=
 \frac{\alpha\beta}{v_n\delta_j^n}\one_G(w)\,dw
\]
is dominated by both
${\sf P}_{\rm loc}(x,\cdot)$ and
${\sf P}_{\rm loc}(y,\cdot)$, and
\[
 \xi(K)
 =
 \frac{\alpha\beta}{v_n\delta_j^n}\vol(G)
 \ge
  \frac{67}{80}\alpha\beta.
\]
Using the fact that two probability measures that dominate a common subprobability
measure of mass $a$ have total-variation distance at most $1-a$, we conclude that
\[
 \dtv\bigl(
 {\sf P}_{\rm loc}(x,\cdot),
 {\sf P}_{\rm loc}(y,\cdot)
 \bigr)
 \le
 1-\frac{67}{80}\alpha\beta.
\]
The result follows with
$\alpha_0:=\frac{67}{80}\alpha\beta$.
\end{proof}

\end{document}
